% Appendix D. Source: Appendix 4.C of Dissertation/Chapter4_Methodology.md.

\chapter{What the 52 Automated Tests Check}
\label{app:tests}
\lhead{Appendix D. \emph{What the 52 Automated Tests Check}}

Section~\ref{sec:verification} reports 52 tests. This appendix says what each one checks, so the count can be verified and any single check found. Every test runs whenever the code changes, and a failing test stops the change rather than reporting the problem afterwards.

\section{The selection model, tests 1 to 35}
\label{sec:testsmodel}

The model of Chapter~\ref{ch:problem}, and the two refinements of Section~\ref{sec:parameters}.

\begin{longtable}{@{}rp{0.86\textwidth}@{}}
  \caption{Tests 1 to 13: the selection model}
  \label{tab:tests1}\\
  \toprule
  & \textbf{The test checks that} \\
  \midrule
  \endfirsthead
  \toprule
  & \textbf{The test checks that} \\
  \midrule
  \endhead
  \bottomrule
  \endfoot
  1  & The demonstration instance loads with 16 trials covering all four pathways \\
  2  & The demonstration durations are short enough to show the on-time mechanism, which the measured durations are not \\
  3  & The solver finds the same plan as checking every possible plan one at a time \\
  4  & The solver agrees with the simpler ranking rule wherever that rule is known to be correct \\
  5  & No plan commits more trials than the limit allows \\
  6  & The model commits fewer trials than the limit when an extra trial would cost a rescue \\
  7  & Ranking on lateness alone rescues no trial, which is why on-time is ranked first \\
  8  & Putting on-time first never improves lateness; it can only cost lateness \\
  9  & A study scoring below the value floor is closed rather than left in the queue \\
  10 & An empty queue returns a plan rather than an error \\
  11 & The capacity check made before solving agrees with what the solver returns \\
  12 & A trial that cannot be planned is returned marked as blocked, not quietly dropped \\
  13 & Solving through the planner interface gives the same plan as using the model directly \\
\end{longtable}

\begin{longtable}{@{}rp{0.86\textwidth}@{}}
  \caption{Tests 14 to 17: the withdrawn trade setting}
  \label{tab:tests2}\\
  \toprule
  & \textbf{The test checks that} \\
  \midrule
  \endfirsthead
  \toprule
  & \textbf{The test checks that} \\
  \midrule
  \endhead
  \bottomrule
  \endfoot
  14 & The withdrawn setting has been removed from the model, not merely set to zero \\
  15 & A request still using it is refused rather than silently given an unchanged plan \\
  16 & Switching the value objective on never costs an on-time trial \\
  17 & Portfolio value rises with the commitment limit, which is why trading a trial away never pays \\
\end{longtable}

\begin{longtable}{@{}rp{0.86\textwidth}@{}}
  \caption{Tests 18 to 23: segment B, the days lost before a trial reaches pharmacy}
  \label{tab:tests3}\\
  \toprule
  & \textbf{The test checks that} \\
  \midrule
  \endfirsthead
  \toprule
  & \textbf{The test checks that} \\
  \midrule
  \endhead
  \bottomrule
  \endfoot
  18 & A trial already in pharmacy's queue does not lose the same days twice \\
  19 & A trial not yet handed over loses the days before hand-over \\
  20 & A trial R\&I still holds is planned on the median hand-over and marked as assumed \\
  21 & A trial with no usable dates is blocked rather than given an invented window \\
  22 & Subtracting the hand-over can cost a trial its rescue, which is the point of the correction \\
  23 & An assumed hand-over is always reported as assumed, never as a recorded date \\
\end{longtable}

\begin{longtable}{@{}rp{0.86\textwidth}@{}}
  \caption{Tests 24 to 35: segment A, the preparation credit}
  \label{tab:tests4}\\
  \toprule
  & \textbf{The test checks that} \\
  \midrule
  \endfirsthead
  \toprule
  & \textbf{The test checks that} \\
  \midrule
  \endhead
  \bottomrule
  \endfoot
  24 & With no credit set, every plan matches the plan produced before the credit existed \\
  25 & With both conditions met, the credit shortens the set-up \\
  26 & Complexity alone earns nothing without the sponsor green light \\
  27 & The green light alone earns nothing without complexity \\
  28 & A Phase 1 study earns the credit without needing aseptic work \\
  29 & Every input to the credit comes from the data request; none is invented \\
  30 & The planner can add a study the rule would miss \\
  31 & The planner can withhold the credit from a study the rule would give it to \\
  32 & The planner cannot waive the sponsor green light \\
  33 & The credit never shortens a study below one day \\
  34 & A negative credit is refused \\
  35 & The credit can turn a missed deadline into a rescue \\
\end{longtable}

\section{The robustness experiment, tests 36 to 40}
\label{sec:testsrobust}

Tests 36 and 37 are the two hand-worked checks of Section~\ref{sec:verification}.

\begin{longtable}{@{}rp{0.86\textwidth}@{}}
  \caption{Tests 36 to 40: the robustness experiment}
  \label{tab:tests5}\\
  \toprule
  & \textbf{The test checks that} \\
  \midrule
  \endfirsthead
  \toprule
  & \textbf{The test checks that} \\
  \midrule
  \endhead
  \bottomrule
  \endfoot
  36 & The four-trial instance of Appendix~\ref{app:notation} gives the same answer worked by hand and computed by the model \\
  37 & The lateness calculation matches equation~(\ref{eq:lateness}) worked by hand \\
  38 & A scenario regenerates exactly from its seed \\
  39 & The comparator commits the longest-waiting trials, as Section~\ref{sec:comparison} describes \\
  40 & Checking every plan one at a time and solving directly return the same plan \\
\end{longtable}

\section{The plan-level uncertainty run, tests 41 to 52}
\label{sec:testsuncertainty}

Tests 41 to 43 fix the run's agreement with the plan it qualifies.

\begin{longtable}{@{}rp{0.86\textwidth}@{}}
  \caption{Tests 41 to 52: the plan-level uncertainty run}
  \label{tab:tests6}\\
  \toprule
  & \textbf{The test checks that} \\
  \midrule
  \endfirsthead
  \toprule
  & \textbf{The test checks that} \\
  \midrule
  \endhead
  \bottomrule
  \endfoot
  41 & With no variation in duration, the run returns the planning model's on-time count exactly \\
  42 & With variation, the promised count is not always achieved, which is the point of the run \\
  43 & Only the set-up spread widens the drawn durations; the approval-time figure cannot \\
  44 & The reported distribution adds to one and never exceeds the number of trials committed \\
  45 & The pessimistic figure never reads higher than the average \\
  46 & The trial most likely to miss its deadline is reported first \\
  47 & A trial the plan cannot rescue is never shown as rescued \\
  48 & A larger plan carries a longer approval wait \\
  49 & A run repeats exactly from the same seed \\
  50 & A different seed gives a different result \\
  51 & An empty plan is refused \\
  52 & A run must have at least one repetition \\
\end{longtable}
